\documentclass[11pt,reqno]{amsart}
\usepackage[T1]{fontenc}
\usepackage{amsmath,amssymb,amsthm,mathtools}
\usepackage[colorlinks=true,linkcolor=blue,citecolor=blue,urlcolor=blue]{hyperref}
\newtheorem{theorem}{Theorem}[section]
\newtheorem{proposition}[theorem]{Proposition}
\newtheorem{lemma}[theorem]{Lemma}
\newtheorem{corollary}[theorem]{Corollary}
\theoremstyle{definition}\newtheorem{definition}[theorem]{Definition}
\theoremstyle{remark}\newtheorem{remark}[theorem]{Remark}
\DeclareMathOperator{\diag}{diag}
\DeclareMathOperator{\tr}{tr}
\DeclareMathOperator{\spec}{spec}
\DeclareMathOperator{\supp}{supp}
\newcommand{\RR}{\mathbb R}
\newcommand{\NN}{\mathbb N}
\newcommand{\ip}[2]{\langle #1,#2\rangle}
\newcommand{\norm}[1]{\lVert #1\rVert}
\title{Symmetric Krawtchouk Matrices:\\A Commuting Difference Operator and a Squeezed Mehler Limit}
\date{Working research note; Lie-algebra interpretation incorporated}
\keywords{Krawtchouk matrix, commuting difference operator, Perron eigenvector, Rayleigh quotient, Mehler kernel, squeezing, Lie algebra}
\begin{document}
\begin{abstract}
The skew-symmetric operator of Feinsilver--Fitzgerald anticommutes with
the symmetric Krawtchouk matrix $S_N$. A sign conjugation turns its
negative square into positive irreducible Jacobi blocks on the parity
sublattices, with the same Perron eigenvectors as the corresponding
blocks of $S_N^2$. Gaussian Rayleigh quotients select a width
$\alpha=\sqrt8$ and yield the asymptotic lower constant
$(\sqrt8-2)/\pi$ for $N\lambda_{\max}(S_N)^2/8^N$.
The central-window kernel is a squeezed Mehler semigroup. We compare
this construction with the finite oscillator and fractional
Fourier--Kravchuk framework of Atakishiyev, Vicent, and Wolf, and use
supplied preliminary notes to express squeezing and imaginary-time
evolution through an $\mathfrak{sl}(2)$ triple. Exact finite identities,
proved lower bounds, exact properties of the candidate limit, and
unproved convergence assertions are separated throughout.
\end{abstract}
\maketitle
\tableofcontents

\section{Conventions, scope, and proof status}
Define the ordinary-generating-function polynomials by
\begin{equation}\label{eq:gen}
(1+v)^{N-j}(1-v)^j=\sum_{n=0}^N K_n(j;N)v^n,
\qquad 0\le j\le N.
\end{equation}
Set
\begin{equation}\label{eq:matrices}
\Phi_N=(K_n(j;N))_{n,j=0}^N,\qquad
B_N=\diag\left(\binom Nj\right)_{j=0}^N,\qquad
S_N=\Phi_NB_N.
\end{equation}
All indices start at zero; binomial coefficients outside their support
are zero. The weighted duality
\[
\binom NjK_n(j;N)=\binom NnK_j(n;N)
\]
makes $S_N$ real symmetric. Throughout, the matrix called $H_N$ below
is a positive finite difference operator, not the oscillator Hamiltonian
$\mathscr H_\omega$ introduced later.

The finite anticommutation, exceptional eigenvector, and spectral pairing
are from Feinsilver--Fitzgerald \cite{FF}. The finite-oscillator and
fractional-transform background is from \cite{AW,AVW}. The supplied
preliminary notes \cite{Notes} motivate our explicit algebraic
interpretation. We make no bibliographic priority claim for the further
consequences developed here.

\subsection*{What is established and what remains open}
We prove the common-Perron-mode reduction, compact-window Rayleigh
limits, an asymptotic lower bound, and exact spectral formulas for the
candidate Mehler integral operator. The predicted equivalence
\begin{equation}\label{eq:predicted-asymptotic}
\lambda_{\max}(S_N)\sim
\frac{\sqrt{\sqrt8-2}}{\sqrt{\pi N}}(\sqrt8)^N
\end{equation}
is \emph{not} asserted as a theorem in this note. A matching upper bound
or suitable global operator convergence remains to be established.
Likewise, the spectrum of the candidate limiting kernel is not, by
itself, a proof of convergence of the finite excited eigenvalues.

\section{The finite skew-symmetric operator}
\subsection{The finite $\mathfrak{sl}(2)$ representation}
On the standard coordinate basis let
\[
R_Ne_n=(N-n)e_{n+1},\qquad L_Ne_n=ne_{n-1},\qquad
\Lambda_Ne_n=(N-2n)e_n.
\]
Endpoint terms are omitted. Direct calculation gives
\begin{equation}\label{eq:finite-sl2}
[L_N,R_N]=\Lambda_N,\qquad
[\Lambda_N,R_N]=-2R_N,\qquad
[\Lambda_N,L_N]=2L_N.
\end{equation}
These are an $\mathfrak{sl}(2)$ representation after a choice of sign
and generator convention. They need not be adjoints in the ordinary
Euclidean coordinate inner product.

The operator in \cite{FF} is
\begin{equation}\label{eq:T-enveloping}
T_N=(\Lambda_N-NI)R_N+(\Lambda_N+NI)L_N.
\end{equation}
Thus $T_N$ is an explicit element of the universal enveloping algebra,
quadratic plus linear in these generators. Equivalently,
\begin{equation}\label{eq:T-action}
\frac12T_Ne_n=n(N-n+1)e_{n-1}-(n+1)(N-n)e_{n+1}.
\end{equation}
For $a_n=2(n+1)(N-n)$, $0\le n<N$, its superdiagonal is $a_n$
and its subdiagonal is $-a_n$.

\begin{proposition}[Feinsilver--Fitzgerald \cite{FF}]
\label{prop:T}
The operator satisfies
\begin{equation}\label{eq:anticommutation}
T_N^T=-T_N,\qquad T_NS_N=-S_NT_N.
\end{equation}
For odd $N$ its kernel is zero. For even $N=2M$, its kernel is spanned
by the vector $z$ given by
\begin{equation}\label{eq:exceptional-vector}
z_{2m}=\binom{2m}{m}\binom{N-2m}{M-m},\qquad z_{2m+1}=0,
\end{equation}
and $S_Nz=\binom N{N/2}z$.
\end{proposition}
Skew symmetry is immediate from \eqref{eq:T-action}. The null-vector
recurrence and anticommutation are proved in \cite{FF}; they can also
be verified from \eqref{eq:gen} and the displayed finite generators.
In particular,
\begin{equation}\label{eq:square-commutation}
[T_N,S_N^2]=0,\qquad [T_N^2,S_N]=0.
\end{equation}
Nonzero images under $T_N$ pair eigenvectors for $\lambda$ and $-\lambda$.
For odd $N$ the characteristic polynomial is even; for even $N$ it is
$(\lambda-\binom N{N/2})$ times an even polynomial, as stated in \cite{FF}.

\section{Positive parity blocks and the common Perron mode}
Let
\begin{equation}\label{eq:sign-conjugation}
\sigma_N=\diag((-1)^{\lfloor n/2\rfloor}),\qquad
H_N=\sigma_N(-T_N^2)\sigma_N,\qquad
G_N=\sigma_NS_N^2\sigma_N.
\end{equation}
Since $\sigma_{n+2}=-\sigma_n$, the nonzero entries of $H_N$ are
\begin{equation}\label{eq:H-entries}
(H_N)_{nn}=a_{n-1}^2+a_n^2,\qquad
(H_N)_{n,n+2}=a_na_{n+1}=(H_N)_{n+2,n},
\end{equation}
with $a_{-1}=a_N=0$. Each parity restriction is a nonnegative,
irreducible symmetric tridiagonal matrix. For a singleton block the
same statements about its positive Perron vector are interpreted directly.

\subsection{A positive formula for $G_N$}
Put $d_i=\binom Ni$. Parseval on the generating functions gives
\begin{equation}\label{eq:bilinear}
\sum_nK_n(i;N)K_n(j;N)
=\frac{4^N}{\pi}\cos\frac{\pi(i-j)}2\,
\mathrm B\left(\frac{i+j+1}{2},\frac{2N-i-j+1}{2}\right).
\end{equation}
Because $S_N$ is symmetric, $(S_N^2)_{ij}$ is $d_id_j$ times this
column inner product. Equal-parity indices satisfy
$\sigma_i\sigma_j\cos(\pi(i-j)/2)=1$. Therefore
\begin{equation}\label{eq:G-entries}
(G_N)_{ij}=\begin{cases}
\displaystyle d_id_j\frac{4^N}{\pi}
\mathrm B\left(\frac{i+j+1}{2},\frac{2N-i-j+1}{2}\right),
&i\equiv j\pmod2,\\[4pt]
0,&i\not\equiv j\pmod2.
\end{cases}
\end{equation}
Each parity block is strictly positive, and $[H_N,G_N]=0$.

\begin{theorem}[Common Perron eigenvectors]\label{thm:common-PF}
On each parity sublattice the Perron eigenvector of $H_N$ is the
Perron eigenvector of $G_N$. For $N\ge1$, the two parity blocks of
$H_N$ have the same largest eigenvalue, and the two parity blocks of
$G_N$ have the same largest eigenvalue. Moreover,
$\lambda_{\max}(S_N)=\norm{S_N}_2$.
\end{theorem}
\begin{proof}
The top eigenspace of each irreducible Jacobi block of $H_N$ is
one-dimensional and has a positive vector $u$. Commutation implies
$G_Nu=cu$. Since the corresponding block of $G_N$ is strictly positive,
this is its Perron eigenpair.

The operator $T_N$ exchanges parity and commutes with $-T_N^2$.
Every positive eigenvalue of $-T_N^2$ therefore has equal multiplicities
on the two parity subspaces: if $-T_N^2v=hv$ with $h>0$, then
$T_Nv\ne0$, and $T_N$ has inverse $-T_N/h$ on that eigenspace.
Consequently the largest $H_N$ eigenvalues agree. Their common
positive eigenspace, before sign conjugation, has dimension two.
The commuting $S_N^2$ acts on each parity line by its respective
Perron eigenvalue; commutation with $T_N$ makes these scalars equal.
Thus this two-dimensional space carries the largest eigenvalue of
$S_N^2$. On it, $T_N$ is invertible and anticommutes with $S_N$,
so the two $S_N$ eigenvalues are opposite. This proves the last assertion.
\end{proof}
The eigenvalues of $H_N$ and $G_N$ are \emph{not} equated by this
argument. It identifies their largest-mode eigenvectors, which is the
useful reduction from a dense matrix to a sparse difference operator.

\section{Gaussian Rayleigh quotients for the sparse operator}
Fix a parity and set
\begin{equation}\label{eq:trial}
y_n=\frac{n-N/2}{\sqrt N},\qquad
f_n^{(\alpha)}=e^{-\alpha y_n^2/2},\qquad \alpha>0.
\end{equation}
The exact sparse Rayleigh quotient is
\begin{equation}\label{eq:RH-exact}
R_{H,N}(\alpha)=
\frac{\sum_n(a_{n-1}^2+a_n^2)f_n^2
+2\sum_n a_na_{n+1}f_nf_{n+2}}{\sum_n f_n^2}.
\end{equation}
Each sum is over the appropriate parity indices and valid edges.

\begin{proposition}[Sparse Gaussian quotient]\label{prop:RH}
For fixed $\alpha>0$,
\begin{equation}\label{eq:RH-asymptotic}
R_{H,N}(\alpha)=N^4+N^3\left(4-\frac\alpha2-\frac4\alpha\right)+O(N^2).
\end{equation}
The $N^3$ coefficient is maximized at $\alpha=\sqrt8$.
\end{proposition}
\begin{proof}
Let $r_n$ be the row sum of $H_N$. The exact energy identity is
\[
\ip f{H_Nf}=\sum_n r_nf_n^2
-\sum_n a_na_{n+1}(f_{n+2}-f_n)^2.
\]
Polynomial expansion about $n=N/2+\sqrt N\,y$ gives
\[
r_n=N^4+N^3(4-8y^2)
+N^2(16y^4-16y^2)+48Ny^2-8N.
\]
The parity mesh in $y$ is $h_N=2/\sqrt N$.
Gaussian Riemann sums yield
$\sum y_n^2f_n^2/\sum f_n^2=1/(2\alpha)+o(1)$,
with exponentially small sampling and endpoint errors. On the central
scale $a_na_{n+1}=N^4/4+O(N^{7/2}|y|+N^3(1+y^2))$;
expanding the edge difference gives the leading energy
$N^3\int(f')^2/\int f^2=N^3\alpha/2$.
Using the edge midpoint, or cancellation of odd terms, bounds the
remaining contribution by $O(N^2)$ for this fixed Gaussian.
Gaussian tails control all polynomial remainder terms.
This proves \eqref{eq:RH-asymptotic}.
Finally $\alpha/2+4/\alpha\ge\sqrt8$, with equality at $\sqrt8$.
\end{proof}
The local differential expression suggested by this calculation is
\begin{equation}\label{eq:local-H}
H_N=N^4+N^3\left(\frac{d^2}{dy^2}+4-8y^2\right)+\cdots.
\end{equation}
Equation \eqref{eq:local-H} is a local scaling expression, not an
assertion of global operator convergence. Proposition~\ref{prop:RH}
does give a genuine trial lower bound with leading terms
$N^4+(4-\sqrt8)N^3$ for the largest eigenvalue of $H_N$.

\section{The central kernel and a rigorous spectral lower bound}
\subsection{Compact-window Stirling asymptotics}
Write $i=N/2+\sqrt N\,y$ and $j=N/2+\sqrt N\,z$.
Uniformly for $y,z$ in a fixed compact set, the binomial and beta
factors in \eqref{eq:G-entries} give, on equal-parity indices,
\begin{equation}\label{eq:G-local}
(G_N)_{ij}=\frac{8^N}{N^{3/2}}\frac{2\sqrt2}{\pi^{3/2}}
 e^{-3(y^2+z^2)/2+yz}\bigl(1+O(N^{-1/2})\bigr).
\end{equation}
Indeed,
\[
\binom Ni=2^N\sqrt{\frac2{\pi N}}e^{-2y^2}
(1+O(N^{-1/2})),
\]
and the beta factor is
\[
2^{-N}\sqrt{\frac{2\pi}N}\,e^{(y+z)^2/2}
(1+O(N^{-1/2})).
\]
The parity spacing $2/\sqrt N$, rather than $1/\sqrt N$, is essential
to the resulting integral-operator constant.

Define the candidate limiting operator on $L^2(\RR,dy)$ by
\begin{equation}\label{eq:limiting-kernel}
(\mathcal Gf)(y)=\frac{\sqrt2}{\pi^{3/2}}
\int_{\RR}e^{-3(y^2+z^2)/2+yz}f(z)\,dz.
\end{equation}
For nonzero continuous compactly supported $f$, sampling $f(y_n)$
on either parity gives the proved Rayleigh limit
\begin{equation}\label{eq:compact-Rayleigh}
\lim_{N\to\infty}\frac N{8^N}
\frac{\ip{f_N}{G_Nf_N}}{\ip{f_N}{f_N}}
=\frac{\ip f{\mathcal Gf}}{\ip f f}.
\end{equation}
This follows from uniform compact asymptotics and Riemann sums; it
does not require control of matrix entries outside that compact window.

\subsection{Optimizing the limiting Gaussian quotient}
For $f_\alpha(y)=e^{-\alpha y^2/2}$, Gaussian integration gives
\begin{equation}\label{eq:limiting-Rayleigh}
\frac{\ip{f_\alpha}{\mathcal Gf_\alpha}}{\ip{f_\alpha}{f_\alpha}}
=\frac{2\sqrt{2\alpha}}{\pi\sqrt{\alpha^2+6\alpha+8}}.
\end{equation}
The maximum over this Gaussian family is attained at
\begin{equation}\label{eq:optimal-parameters}
\alpha=\sqrt8,\qquad c=\frac{\sqrt8-2}{\pi}.
\end{equation}
For example, differentiating the logarithm of the quotient leaves
$\alpha^2=8$. This is the same width selected by the sparse quotient.

\begin{theorem}[Asymptotic Rayleigh lower bound]\label{thm:lower-bound}
\begin{equation}\label{eq:lower-bound}
\boxed{\liminf_{N\to\infty}
\frac{\pi N\lambda_{\max}(S_N)^2}{8^N}\ge\sqrt8-2.}
\end{equation}
\end{theorem}
\begin{proof}
Choose compactly supported continuous cutoffs of $f_{\sqrt8}$.
For each fixed cutoff, apply Rayleigh's principle and
\eqref{eq:compact-Rayleigh}. Then enlarge the cutoff.
The operator \eqref{eq:limiting-kernel} is Hilbert--Schmidt because
its squared kernel has a strictly negative definite quadratic exponent;
its Rayleigh quotient is continuous under $L^2$ approximation.
Thus the cutoff quotients tend to $c$. The largest eigenvalue of each
$G_N$ block is $\lambda_{\max}(S_N)^2$ by
Theorem~\ref{thm:common-PF}. This proves the bound.
\end{proof}
No step of this proof replaces a missing global upper estimate by a
formal Gaussian tail assertion for the finite matrices.

\section{The exact Mehler structure of the candidate limit}
For $\alpha>0$ let
\[
\psi_r^{(\alpha)}(y)=
\frac{\alpha^{1/4}}{\pi^{1/4}\sqrt{2^r r!}}
H_r(\sqrt\alpha\,y)e^{-\alpha y^2/2},
\]
where $H_r$ are physicists' Hermite polynomials. Define
\[
\mathcal N_\alpha=
\frac1{2\alpha}\left(-\frac{d^2}{dy^2}+\alpha^2y^2\right)-\frac12.
\]
Its normalized eigenfunctions are $\psi_r^{(\alpha)}$ with eigenvalues $r$.

\begin{proposition}[Mehler identification]\label{prop:Mehler}
With $\alpha=\sqrt8$, $q=3-\sqrt8$, and $c=(\sqrt8-2)/\pi$,
\begin{equation}\label{eq:Mehler-expansion}
\mathcal G(y,z)=c\sum_{r=0}^\infty
q^r\psi_r^{(\alpha)}(y)\psi_r^{(\alpha)}(z).
\end{equation}
Consequently the eigenvalues of $\mathcal G$ are $cq^r$, and
\begin{equation}\label{eq:Mehler-semigroup}
\mathcal G=c\exp(-\tau\mathcal N_\alpha),\qquad
\tau=-\log q=2\log(1+\sqrt2).
\end{equation}
\end{proposition}
\begin{proof}
The normalized Hermite Mehler sum is
\[
\sum_{r\ge0}q^r\psi_r^{(\alpha)}(y)\psi_r^{(\alpha)}(z)
=\sqrt{\frac\alpha\pi}\frac1{\sqrt{1-q^2}}
\exp\left[-\frac{\alpha(1+q^2)}{2(1-q^2)}(y^2+z^2)
+\frac{2\alpha q}{1-q^2}yz\right].
\]
The stated parameters give diagonal coefficient $3/2$, mixed
coefficient $1$, and, after multiplication by $c$, amplitude
$\sqrt2/\pi^{3/2}$. This is \eqref{eq:limiting-kernel}.
\end{proof}
The Gaussian Rayleigh vector is thus the \emph{exact} ground state of
the candidate integral operator. Excited Hermite modes have geometric
multipliers $q^r$. These statements about $\mathcal G$ are exact, while
their transfer to the finite spectrum is conditional on a suitable
convergence theorem.

% ALGEBRA-SECTION-START
\section{The Lie-algebra interpretation: squeezing and imaginary time}
\label{sec:Lie-algebra}
The supplied preliminary notes \cite{Notes} provide a useful algebraic
interpretation. We spell it out with explicit coordinate and domain
conventions to distinguish the finite representation from the oscillator
representation.

\subsection{The quadratic oscillator realization}
On the Schwartz space $\mathcal S(\RR)$, let $D=d/dx$ and set
\begin{equation}\label{eq:oscillator-sl2}
\mathsf R=\frac{x^2}{2},\qquad
\mathsf L=\frac{D^2}{2},\qquad
\rho=xD+\frac12.
\end{equation}
The Heisenberg relation $[D,x]=1$ gives
\begin{equation}\label{eq:oscillator-sl2-brackets}
[\mathsf L,\mathsf R]=\rho,\qquad
[\mathsf L,\rho]=2\mathsf L,\qquad
[\rho,\mathsf R]=2\mathsf R.
\end{equation}
Thus this is an $\mathfrak{sl}(2)$ realization; a sign change of a
lowering generator gives another common convention.
The frequency-$\omega$ Hamiltonian and number operator are
\begin{equation}\label{eq:omega-oscillator}
\mathscr H_\omega=\frac12(-D^2+\omega^2x^2),\qquad
\mathcal N_\omega=\frac{\mathscr H_\omega}{\omega}-\frac12
=\omega\mathsf R-\omega^{-1}\mathsf L-\frac12.
\end{equation}
The natural Hilbert space for the wavefunctions and the unitary Fourier
transform is $L^2(\RR,dx)$. A Gaussian-weighted polynomial realization
requires conjugating the operators; it is not the same realization
with its measure changed silently.

\subsection{Squeezing is the exponential of $\rho$}
The closure of $\rho$ generates the unitary dilation group
\begin{equation}\label{eq:dilation}
V_r=e^{r\rho},\qquad (V_rf)(x)=e^{r/2}f(e^rx),\qquad r\in\RR.
\end{equation}
Direct calculation gives
\[
V_rxV_r^{-1}=e^rx,\qquad V_rDV_r^{-1}=e^{-r}D.
\]
For $e^{2r}=\omega$ this implies
\begin{equation}\label{eq:squeezed-conjugacy}
V_r\mathscr H_1V_r^{-1}=\frac{\mathscr H_\omega}{\omega},\qquad
\mathcal N_\omega=V_r\mathcal N_1V_r^{-1}.
\end{equation}
In particular,
\[
V_r(\pi^{-1/4}e^{-x^2/2})
=(\omega/\pi)^{1/4}e^{-\omega x^2/2}.
\]
Thus the changed width of our Perron profile is an actual dilation
in the oscillator representation, not only an analogy with squeezing.

\subsection{The limiting kernel as an algebraic semigroup}
Using \eqref{eq:omega-oscillator} in the variable $y$, the exact
Mehler operator can be written
\begin{equation}\label{eq:Lie-semigroup}
\boxed{
\mathcal G=c\,e^{\tau/2}
\exp\left[\tau(\alpha^{-1}\mathsf L-\alpha\mathsf R)\right],
\qquad \alpha=\sqrt8.
}
\end{equation}
Equivalently,
\[
\mathcal G=cV_{\frac12\log\alpha}
 e^{-\tau\mathcal N_1}V_{\frac12\log\alpha}^{-1}.
\]
For positive real $\tau$, this is a bounded imaginary-time oscillator
semigroup, rather than a unitary real-time rotation. Algebraically it
is the semigroup continuation of the same quadratic oscillator
construction used for fractional Fourier transformations.

With the positive-sign Fourier convention of the supplied notes,
\[
(\mathcal F_+f)(\xi)=\frac1{\sqrt{2\pi}}\int e^{i\xi x}f(x)\,dx,
\qquad
\mathcal F_+=e^{i\pi\mathcal N_1/2}.
\]
This contrasts with the negative-sign convention used in much of the
optics literature. The notes derive this identity from
$\psi_n=(x-D)^ne^{-x^2/2}$ and
$\mathscr H_1\psi_n=(n+1/2)\psi_n$.
In their generating-function calculation the intermediate exponent
is $-\xi^2/2+2i\xi t+t^2$; the factor $2$ in the mixed term is needed.
We use normalized Hermite wavefunctions and the unweighted $L^2$ space
when passing from the algebraic calculation to unitary operators.

\subsection{What this says, and does not say, about the finite algebra}
Equation \eqref{eq:T-enveloping} places the finite $T_N$ in the universal
enveloping algebra of \eqref{eq:finite-sl2}. Equation
\eqref{eq:Lie-semigroup} places the candidate limiting kernel in the
oscillator semigroup associated with \eqref{eq:oscillator-sl2-brackets}.
These are two different representations. We have not exhibited an
exact finite-$N$ Lie-group conjugacy between them, and $T_N$ is not
simply the standard normalized finite-oscillator generator.
The centered scaling of $H_N$, its Perron localization, and the
weighted kernel asymptotics are the analytic bridge that must be
controlled. The algebraic interpretation organizes that bridge but
does not replace its estimates.

The even and odd Hermite subspaces are each preserved by
$\mathsf R,\mathsf L,\rho$. This oscillator parity should not be
confused with the even/odd \emph{lattice-index} blocks of the finite
matrices. Each finite lattice block samples the entire real line in
the central scaling and can support both even and odd continuum modes.

Finally, the notes ask about an intertwiner with convolution.
Identifying a finite mode basis with the first $N+1$ Hermite modes
intertwines diagonal fractional phases on that subspace. It does not
by itself intertwine ordinary convolution with a finite Kravchuk
convolution, or identify the physical position operators. We leave
that stronger question separate from the spectral argument.
% ALGEBRA-SECTION-END

\section{Comparison with Atakishiyev--Vicent--Wolf}
\label{sec:AVW}
The finite-oscillator interpretation and its fractional Fourier
transforms are established background \cite{AW,AVW}, not novelty
claims of this note. In particular, \cite{AVW}, equations (21)--(28),
organizes the finite oscillator through angular-momentum generators;
equations (33)--(35) give normalized Kravchuk functions and their
oscillator contraction; equation (37) gives the fractional-transform
spectral factorization.

\subsection{A normalization dictionary}
Their variables translate as $N=2\ell$, $m=j-N/2$. Their polynomial
normalization in equation (34), denoted here only in this paragraph
by $k_n^{\mathrm{AVW}}$, satisfies
\[
k_n^{\mathrm{AVW}}(j;N)=(-1)^n2^{-n}K_n(j;N).
\]
This is \emph{not} the factorially normalized centered polynomial that
is also sometimes denoted $k_n$ in other work.
Their normalized functions are consequently
\begin{equation}\label{eq:AVW-functions}
\phi_n(j-N/2)=(-1)^n2^{-N/2}
\sqrt{\frac{\binom Nj}{\binom Nn}}K_n(j;N).
\end{equation}
Set
\begin{equation}\label{eq:U}
U_N=2^{-N/2}B_N^{-1/2}\Phi_NB_N^{1/2}.
\end{equation}
Duality and orthogonality give $U_N^T=U_N$, $U_N^TU_N=I$.
Up to the mode-dependent sign, \eqref{eq:AVW-functions} is its row
basis. The spectral construction can therefore be written in our
convention as
\[
\mathcal F_{N,\theta}
=U_N\diag(e^{-in\theta})_{n=0}^NU_N^T.
\]
The paper's fractional parameter rescales $\theta$; Fourier signs and
global phases must be fixed before equating kernels literally.
The mode signs in \eqref{eq:AVW-functions} cancel in this spectral sum.

For our spectral problem the essential relation is
\begin{equation}\label{eq:congruence}
S_N=2^{N/2}B_N^{1/2}U_NB_N^{1/2}.
\end{equation}
This is a congruence, not a similarity. The unitary transform's simple
spectral structure therefore does not directly determine the growing
eigenvalues of $S_N$.

\subsection{Standard versus squeezed oscillator contraction}
The standard finite ground state is
$2^{-N/2}\sqrt{\binom Nj}$, with central shape $e^{-y^2}$.
The paper's continuum coordinate is
\[
x=\frac m{\sqrt\ell}=\sqrt2\,y.
\]
For unit continuum normalization the discrete functions acquire the
factor $\ell^{1/4}$, the inverse square root of their lattice spacing.
The fixed-degree contraction, with sign conventions consistently
chosen, gives the usual normalized Hermite functions. In particular
the standard ground state has shape $e^{-x^2/2}$.

Our Rayleigh profile is instead
\[
e^{-\sqrt2 y^2}=e^{-\sqrt2 x^2/2}.
\]
It is the frequency-$\sqrt2$ squeezed ground state in the paper's
coordinate, with dilation parameter
\begin{equation}\label{eq:physical-squeeze}
r=\frac14\log2.
\end{equation}
Its width is smaller than the standard ground-state width by
$2^{-1/4}$. The fractional-transform framework supplies the natural
basis and group interpretation; the binomial congruence in
\eqref{eq:congruence} and the commuting $T_N$ construction lead to
this particular squeezing and positive limiting kernel.

\section{Numerical evidence and the remaining proof program}
\subsection{Finite Gaussian Rayleigh calculations}
The following values use the even-parity Gaussian trial
$\alpha=\sqrt8$, without finite-$N$ width optimization. By
Theorem~\ref{thm:common-PF}, $\sqrt{R_{G,N}}$ is a trial lower bound
for $\lambda_{\max}(S_N)$ in exact arithmetic. The displayed values
are floating-point evaluations, not certified numerical enclosures.
\begin{center}\small
\begin{tabular}{r|r|r}
$N$&$\sqrt{R_{G,N}(\sqrt8)}$&Numerical $\lambda_{\max}(S_N)$\\\hline
4&15.91790493&16.00000000\\
10&5254.861928&5256.904661\\
20&$1.225209116\times10^8$&$1.225344027\times10^8$\\
50&$2.736712075\times10^{21}$&$2.736758628\times10^{21}$\\
100&$7.319951192\times10^{43}$&$7.319981980\times10^{43}$
\end{tabular}
\end{center}
The sparse and dense Perron vectors agree to numerical precision,
as predicted by the exact common-mode theorem. Numerical evidence
supports \eqref{eq:predicted-asymptotic} but is not its proof.

\subsection{A route to a matching upper bound}
One option is to embed each parity block into $L^2(\RR)$ by step
functions of width $2/\sqrt N$. The kernel corresponding to
$NG_N/8^N$ then has cell values $(\sqrt N/2)NG_{ij}/8^N$.
Compact-window convergence is \eqref{eq:G-local}; a uniform integrable
bound on the squared kernels would imply Hilbert--Schmidt convergence
and hence operator-norm convergence to $\mathcal G$.
Such a bound must control central, moderate-deviation, and edge
regions. It has not been supplied here.

An alternative is to use the common Perron vector, obtain global
localization and a spectral-gap estimate for the Jacobi operator
$H_N$, and prove sufficient convergence of that vector to the squeezed
Gaussian. Transferring this convergence to the dense quadratic form
still requires an adequate bound on the scaled $G_N$; a local
oscillator expansion alone is insufficient.

Once either route is completed, $\norm{\mathcal G}=c$ gives the missing
upper bound and \eqref{eq:predicted-asymptotic}. Stronger compact
operator convergence would also make the geometric excited-mode
predictions accessible. The finite Lie-algebra structure and the
oscillator semigroup clarify the expected result, but the global
analytic estimates remain the principal unfinished task.

\begin{thebibliography}{9}
\bibitem{FF}
Philip Feinsilver and Robert Fitzgerald,
\emph{The spectrum of symmetric Krawtchouk matrices},
Linear Algebra and its Applications (1996),
\href{https://doi.org/10.1016/0024-3795(94)00123-5}
{doi:10.1016/0024-3795(94)00123-5}.
The supplied source \texttt{kravspec.tex} is used for the operator
$T_N$, its kernel, and the spectral pairing.

\bibitem{AW}
N.~M. Atakishiyev and K.~B. Wolf,
\emph{Fractional Fourier--Kravchuk transform},
Journal of the Optical Society of America A \textbf{14} (1997),
1467--1477,
\href{https://doi.org/10.1364/JOSAA.14.001467}
{doi:10.1364/JOSAA.14.001467}.
Cited as the antecedent of the construction analyzed in \cite{AVW}.

\bibitem{AVW}
N.~M. Atakishiyev, L.~E. Vicent, and K.~B. Wolf,
\emph{Continuous vs.\ discrete fractional Fourier transforms},
Journal of Computational and Applied Mathematics \textbf{107}
(1999), 73--95,
\href{https://doi.org/10.1016/S0377-0427(99)00082-5}
{doi:10.1016/S0377-0427(99)00082-5}.
\href{https://www.fis.unam.mx/~bwolf/Articles/106.pdf}{Author-hosted PDF}.

\bibitem{Notes}
Supplied preliminary notes, dated 28 April 2011,
\texttt{FeinNotes11-2.PDF}, three pages, unpublished.
Sections on Hermite functions, $\mathfrak{sl}(2)$ operators, the
Kravchuk connection, and a question about convolution intertwiners.
\end{thebibliography}
\end{document}
